\worksheettitle{Домашняя работа: максимум и минимум}{\modulemark{M06-H} Самостоятельное закрепление}
\studentnameblock

\begin{notebox}[Перед началом]
  Каждую данную цифру используйте ровно один раз. Для максимума располагайте
  цифры по убыванию. Для минимума сначала выбирайте наименьшую ненулевую
  цифру, затем ставьте нули и остальные цифры по возрастанию.
\end{notebox}

\begin{practicebox}[1. Составьте числа]
  \begin{center}
  \renewcommand{\arraystretch}{1.55}
  \begin{tabular}{c|c|c}
    \toprule
    цифры & максимум & минимум \\
    \midrule
    \(1,\ 4,\ 6,\ 9\) & \answerblank[27mm] & \answerblank[27mm] \\
    \(0,\ 3,\ 5,\ 8\) & \answerblank[27mm] & \answerblank[27mm] \\
    \(0,\ 0,\ 3,\ 8\) & \answerblank[27mm] & \answerblank[27mm] \\
    \(0,\ 2,\ 2,\ 7\) & \answerblank[27mm] & \answerblank[27mm] \\
    \bottomrule
  \end{tabular}
  \end{center}
\end{practicebox}

\begin{practicebox}[2. Сравните по первому различающемуся разряду]
  \begin{enumerate}[label=\textbf{\arabic*.}]
    \item \(5038\) и \(5308\): меньше \answerblank[25mm]; первый
      различающийся разряд \textemdash{} \answerblank[25mm], потому что
      \answerblank[30mm].
    \item \(8610\) и \(8601\): больше \answerblank[25mm]; первый
      различающийся разряд \textemdash{} \answerblank[25mm], потому что
      \answerblank[30mm].
  \end{enumerate}
\end{practicebox}

\begin{warningbox}[3. Исправьте чужое решение]
  Для цифр \(0,\ 1,\ 5,\ 8\) минимум записали \(1508\). Правильный ответ:
  \answerblank[34mm]. Объясните ошибку сравнением двух чисел.
  \writinglines{2}
\end{warningbox}

\newpage
\begin{practicebox}[4. Сформулируйте правило]
  Дополните предложения.
  \begin{enumerate}[label=\textbf{\arabic*.}]
    \item В максимуме цифры расположены по \answerblank[28mm].
    \item В минимуме первая цифра не может быть \answerblank[24mm].
    \item После первой ненулевой цифры в минимуме ставят \answerblank[28mm].
  \end{enumerate}
\end{practicebox}

\begin{practicebox}[5. Несколько нулей и повторов]
  \begin{enumerate}[label=\textbf{\arabic*.}]
    \item Из цифр \(0,\ 0,\ 4,\ 4,\ 7\) составьте максимум и минимум:
      \answerblank[34mm], \answerblank[34mm].
    \item Объясните положение двух нулей в каждом числе.
      \writinglines{2}
  \end{enumerate}
\end{practicebox}

\begin{fillbox}[6. Восстановите набор]
  Максимум равен \(8640\). Какие цифры были даны? \answerblank[48mm].
  Какой минимум можно составить из этих цифр? \answerblank[36mm].
\end{fillbox}

\begin{warningbox}[7. Найдите две разные ошибки]
  Для цифр \(0,\ 2,\ 6,\ 9\) записали: максимум \(9622\), минимум \(0269\).
  Исправьте оба ответа и назовите разные причины ошибок.
  \writinglines{2}
\end{warningbox}

\begin{checkpointbox}[8. Итог]
  Из цифр \(0,\ 0,\ 2,\ 7,\ 9\) составьте:
  \begin{enumerate}[label=\textbf{\arabic*.}]
    \item максимум: \answerblank[38mm];
    \item минимум: \answerblank[38mm];
    \item объясните оба ответа через порядок цифр и первый различающийся
      разряд.
      \writinglines{1}
  \end{enumerate}
\end{checkpointbox}

\begin{reflectionbox}[Перед сдачей]
  \choicebox\ Все цифры использованы. \quad
  \choicebox\ Ноль не стоит первым. \quad
  \choicebox\ Есть объяснение.
\end{reflectionbox}
